\documentclass[10pt,a4paper]{amsart}
\usepackage[margin=19mm]{geometry}
\usepackage{amsmath,amssymb,mathtools}
\usepackage[scr=rsfs,cal=euler]{mathalfa}
\usepackage{xfrac}
\usepackage[unicode,hidelinks,bookmarks=false]{hyperref}
\newcommand{\ie}{i.e.\ }
\newcommand{\cf}{cf.\ }
\newcommand{\resp}{respectively\ }
\newcommand{\Subsection}{\subsection}
\providecommand{\nobreakdash}{\nobreak\hbox{-}}
\setlength{\emergencystretch}{2em}
\multlinegap0pt
\usepackage{amssymb}
\usepackage[shortlabels]{enumitem}
\usepackage{mathtools}
\newtheorem{lemma}{Lemma}
\newcommand{\R}{\mathbb{R}}
\title{Holomorphic functions on geometrically finite quotients of the ball}
\author{William Sarem}
\date{Source: arXiv:2411.16620v2 (CC BY 4.0)}
\begin{document}
\maketitle

\noindent\emph{Source and licence.} Holomorphic functions on geometrically finite quotients of the ball. Version arXiv:2411.16620v2 (CC BY 4.0), distributed by arXiv under the Creative Commons Attribution 4.0 International licence, http://creativecommons.org/licenses/by/4.0/, as recorded in the arXiv licence metadata for this exact version and shown on the abstract page https://arxiv.org/abs/2411.16620v2. Full licence notice: \texttt{licenses/arxiv-2411.16620-CC-BY-4.0.txt}.\par\smallskip

\noindent\textit{Editorial scope.} Counted unit: the article's lemma lem-cvx-psh (source0.tex lines 268--270). The article's complete original proof of it is delivered with it (source0.tex lines 272--274, one \texttt{proof} environment of 3 lines). No mathematics is written, repaired or re-derived: the statement and the proof are the article's own lines, in the article's own notation, and no retained line is substituted or re-flowed.  No further source-local context is needed: every cross-reference of every retained line resolves inside the two delivered spans.  Nothing was omitted from any retained span, so the package carries no omission record.  No retained line contains a citation, so the wrapper carries no bibliography and no bibliography entry.  No retained line contains a figure environment, an image-inclusion command or drawing code, so no image is delivered and the package declares no asset dependency.  Editorial formatting only, in addition: an AMS \texttt{amsart} preamble that restates the article's own declarations and macros used by the retained lines, namely the environments \texttt{lemma} on separate counters, together with the standard packages those lines need.  The retained environments print in this document as Lemma 1 in order of appearance, whereas the article prints the same items with the numbering of its own full document.  The complete native source of the article as distributed in the arXiv e-print for this exact version is bundled unchanged as \texttt{sources/169080/source0.tex}.
\begin{lemma}\label{lem-cvx-psh}
	Let $(M,\omega)$ be a Kähler manifold. Assume that there exists a continuous function $\phi:M\to \R$ which is convex on $M$, and strictly convex on $M\setminus \phi^{-1}(0)$. Then the compact connected subvarieties of positive dimension of $M$ are included in the level set $\phi^{-1}(0)$. If moreover $\phi$ is an exhaustion, then $M$ is holomorphically convex.
\end{lemma}

\begin{proof}
	Let $\phi:M\to \R$ be a continuous function which is convex on $M$, and strictly convex on $M\setminus \phi^{-1}(0)$. Then $\phi$ is plurisubharmonic on $M$ and strictly plurisubharmonic on $M\setminus \phi^{-1}(0)$.  Moreover for any connected compact subvariety $A$ of $M$, the function $\phi\vert_{A}$ is constant by the maximum principle, so $A$ is included in $\phi^{-1}(0)$ or in $M\setminus \phi^{-1}(0)$. Notice that $A$ cannot be contained in $M\setminus \phi^{-1}(0)$ because in that case $\phi\vert_{A}$ would be constant and strictly plurisubharmonic. Thus $A\subset \phi^{-1}(0)$. If moreover $\phi$ is an exhaustion, then $M$ is holomorphically convex by Grauert's theorem.
\end{proof}

\end{document}
